\clearpage
\pdfbookmark[1]{Fourier. 25. Februar 1833.}{anzeige_fourier_128_130}
\SourceLeaf{128}
\begin{center}
\rule{0.55\linewidth}{0.4pt}\par Göttingische gelehrte Anzeigen.\quad 25. Februar 1833.\par \rule{0.55\linewidth}{0.4pt}
\end{center}
\textit{Analysis bestimmter Gleichungen, von} \textsc{Fourier}.\quad\textit{Erster Teil.}\quad Paris 1831. \textit{Bei Firmin Didot frères.}\quad (XXIV und 258 Seiten im Quartformat.)

Nach dem Plan des Verfassers sollte dieses Werk aus zwei Abteilungen bestehen. Die erste sollte neben einem allgemeinen Überblick über das Ganze die beiden ersten Abschnitte enthalten, die zweite Abteilung die fünf übrigen Abschnitte. Der Tod entriss den Verfasser den Wissenschaften, als der Druck gerade begonnen hatte. Die gesamte erste Abteilung, die wir hier durch die Bemühungen von Herrn \textsc{Navier} erhalten, fand sich jedoch nahezu vollständig ausgearbeitet vor, während der Rest von der Vollendung noch weit entfernt war. Dies ist umso bedauerlicher, als aus dem allgemeinen Überblick hervorgeht, dass dafür manche interessante Untersuchungen vorgesehen waren. Wir müssen hier übergehen, was wir hätten erwarten können, wenn dem Verfasser ein längeres Leben vergönnt gewesen wäre, und uns auf den Bericht über das beschränken, was wir tatsächlich erhalten haben.

Das Werk soll die numerische Lösung bestimmter algebraischer Gleichungen in einen völlig sicheren und möglichst einfachen Ablauf bringen. Insbesondere soll die Anzahl der reellen und imaginären Wurzeln mit vollständiger Bestimmtheit ermittelt werden. Die reellen Wurzeln sollen einzeln in feste Grenzen eingeschlossen und die schrittweise Annäherung an ihre Werte, die sich zu beliebiger Genauigkeit führen lässt, ganz methodisch geordnet werden. Die Hauptgrundlage des Ganzen bildet ein eigener Satz des Verfassers,
\SourceLeaf{129}
der als eine glückliche Verallgemeinerung des Satzes von \textsc{Descartes} (gewöhnlich nach \textsc{Harriot} benannt) anzusehen ist und im Wesentlichen Folgendes besagt. Es sei $X$ ein geordnetes Polynom in $x$ mit ausschließlich reellen Zahlenkoeffizienten. Ferner seien $a,a'$ zwei beliebige bestimmte, ungleiche reelle Größen, wobei $a'-a$ positiv ist. Durch die Substitutionen $x=y+a,\ x=y+a'$ gehe $X$ in die ebenfalls geordneten Funktionen $Y,Y'$ über, in deren Koeffizienten der Reihe nach $g,g'$ Vorzeichenfolgen vorkommen. Der Ausnahmefall, dass ein oder mehrere Koeffizienten verschwinden, wird hier der Kürze wegen beiseitegelassen. Unter diesen Voraussetzungen kann die Gleichung $X=0$ innerhalb der Grenzen $x=a,\ x=a'$ nicht mehr als $g'-g$ reelle Wurzeln enthalten. Genauer gilt: Ist die Anzahl der reellen Wurzeln zwischen jenen Grenzen $=\lambda$, so ist $g'-g-\lambda$ stets entweder $=0$ oder eine gerade positive Zahl. Ist demnach die Differenz $g'-g$, die niemals negativ sein kann, $=0$, so gibt es zwischen jenen Grenzen überhaupt keine reelle Wurzel. Gilt $g'-g=1$, so enthalten die Grenzen eine reelle Wurzel und nicht mehr. Ist schließlich $g'-g=2$, so bleibt vorerst ungewiss, ob zwei reelle Wurzeln zwischen den Grenzen liegen oder überhaupt keine. Statt zweier Grenzen $a,a'$ lässt sich eine größere Zahl auf ganz ähnliche Weise behandeln und so wählen, dass erstens der kleinsten eine Funktion $Y$ mit ausschließlich Vorzeichenwechseln entspricht, der größten eine mit ausschließlich Vorzeichenfolgen. Werden also die entsprechenden Anzahlen der Vorzeichenfolgen der Reihe nach mit $g,g',g'',g'''$ usw. bezeichnet, so ist die erste dieser Zahlen $=0$ und die letzte gleich dem Grad der Gleichung. Zweitens sollen sämtliche einzelnen Differenzen $g'-g,\ g''-g',\ g'''-g''$ usw. nur entweder $=1$ oder $=2$ sein. Dadurch werden sämtliche reellen Wurzeln so zwischen Grenzen eingeschlossen, dass in jedem einzelnen Intervall entweder eine liegen muss oder zwei liegen können. Wie im letzteren Fall ganz methodisch, nötigenfalls durch weitere Verengung der Grenzen, entschieden wird, ob die zwei reellen Wurzeln tatsächlich vorhanden sind oder fehlen, lässt sich hier aus Platzgründen nicht näher ausführen. Wir bemerken nur, dass es immer dann, wenn dieser letzte Fall eintritt, innerhalb solcher Grenzen einen Zwischenwert gibt, bei dem in der Funktion $Y$ ein Koeffizient vor dem letzten verschwindet, während der vorhergehende und der folgende dasselbe Vorzeichen haben. \textsc{Fourier} nennt solche Stellen kritisch. Jede solche kritische
\SourceLeaf{130}
Stelle bedingt also das Fehlen von zwei reellen Wurzeln. Wenn \textsc{Fourier} sich aber zugleich so ausdrückt, dass jedes Mal zwei solche ausfallenden reellen Wurzeln \textit{imaginär} werden, können wir diesen Ausdruck nicht ganz billigen, weil er leicht zu einer Fehlinterpretation führen könnte. Tatsächlich enthält die Gleichung $X=0$ zwar \textit{zusammengezählt} genau so viele Paare imaginärer Wurzeln, wie solche Ausfälle oder kritischen Stellen vorkommen. Die Werte aller imaginären Wurzeln sind jedoch an sich ebenso bestimmte Größen wie die reellen. Jener Ausdruck kann daher leicht so verstanden werden, als gehöre zu jeder \textit{bestimmten} Lücke ein \textit{bestimmtes} Paar imaginärer Wurzeln. Dies ist nicht nur von \textsc{Fourier} \textit{nicht} nachgewiesen, sondern muss zweifelhaft bleiben, solange tiefergehende Untersuchungen diesen interessanten Punkt noch nicht geklärt haben. Damit soll im Übrigen nicht gesagt werden, dass \textsc{Fourier} selbst den Ausdruck so verstanden habe. Wir möchten eher das Gegenteil annehmen und fast vermuten, dass er über die Existenz oder Nichtexistenz eines solchen bestimmten Zusammenhangs im Ungewissen blieb und absichtlich einer offenen Erklärung zu dem verfänglichen Ausdruck auswich. Überhaupt ist \textsc{Fourier} in diesem Werk auf das Wesen und die Berechnung der imaginären Wurzeln gar nicht eingegangen. Daher bleibt noch ein weites Feld zu bearbeiten.\par
Einem so gewandten Mathematiker, wie \textsc{Fourier} es war, konnte es, einmal im Besitz jenes schönen Satzes --- und nach den von Herrn \textsc{Navier} mitgeteilten Notizen war er dies schon seit sehr langer Zeit ---, nicht schwerfallen, darauf die Organisation der Technik zur numerischen Lösung von Gleichungen zu gründen. Die Entwicklung ist mit großer Vollständigkeit und Ausführlichkeit dargestellt. Geübtere Leser würden vielleicht eine etwas gedrängtere Darstellung und die Streichung mancher Wiederholungen vorziehen. Den weniger Geübten werden die vielen gut gewählten und ausführlich behandelten Beispiele willkommen sein. Jedenfalls sichert dieses Werk dem Namen \textsc{Fourier} auch in diesem Teil der Mathematik einen ehrenvollen Platz, den er in anderen schon längst einnimmt.\par
\par\noindent\hfill\rule{0.30\linewidth}{0.4pt}\hfill\mbox{}\par
